% =========================================================================
% ALGORITMOS DE PRÉ-FLUXO E ALTURAS (PUSH-RELABEL)
% =========================================================================
\chapter{Algoritmos de Pré-fluxo e Alturas (Abordagem Dual)}\label{sec:prefluxo_alturas}
\label{sec:prefluxo}

A família de algoritmos de Goldberg e Tarjan~\cite{goldberg1988new} baseia-se no conceito de pré-fluxo e na manutenção de uma função de alturas, caracterizando uma abordagem dual para o Problema do Fluxo Máximo.

Diferente do paradigma primal, o método de pré-fluxo relaxa a restrição de conservação nos vértices intermediários durante a execução, mantendo como invariante a inexistência de caminhos aumentantes entre fonte e sumidouro na rede residual $D_f$. Inicializa-se a rede saturando todos os arcos de saída da fonte $s$, com $h(s) = |V|$ e $h(t) = 0$. Sob a condição $h(u) \le h(v) + 1$ para todo arco residual $(u,v) \in A_f$, não existem caminhos de $s$ a $t$ em $D_f$. O algoritmo transfere excessos localmente através de operações de empurrar (\textit{push}) e ajuste de altura (\textit{relabel}) até que todos os excessos cheguem ao sumidouro $t$ ou retornem à fonte $s$. Quando todos os nós intermediários têm excesso nulo, a conservação de fluxo é restabelecida e o fluxo obtido é máximo.

Formalmente, um pré-fluxo $f$ satisfaz as restrições de capacidade ($0 \le f(u,v) \le c(u,v)$) e admite excesso não negativo nos vértices intermediários $v \in V \setminus \{s, t\}$:
\begin{equation}
    e(v) = \sum_{(u,v) \in A} f(u,v) - \sum_{(v,w) \in A} f(v,w) \ge 0
\end{equation}

Um vértice intermediário com $e(v) > 0$ é denominado ativo. O algoritmo movimenta esse excesso até o sumidouro $t$ ou, quando não houver arcos residuais descendentes, de volta à fonte $s$~\cite{cook1998combinatorial}.

Quando um vértice ativo $u$ tem $h(u) \ge |V|$, seus arcos admissíveis apontam em direção à fonte, viabilizando o retorno do excesso via arcos residuais reversos criados na saturação inicial.

A função de alturas $h: V \rightarrow \mathbb{N}$ é válida para um pré-fluxo se satisfaz:
\begin{enumerate}
    \item $h(t) = 0$;
    \item $h(s) = |V|$;
    \item $h(u) \le h(v) + 1$, para todo $(u,v) \in A_f$.
\end{enumerate}

A terceira condição garante que nenhum caminho direcionado de $s$ a $t$ existe em $D_f$. O algoritmo encerra quando $e(v) = 0$ para todo $v \in V \setminus \{s, t\}$, momento em que o pré-fluxo é um fluxo viável e máximo.

\begin{figure}[!ht]
\centering
\begin{tikzpicture}[thick]
  \node[source node] (s) at (0, 0) {$s$};
  \node[active node] (a) at (2, 1) {$a$};
  \node[flow node] (b) at (2, -1) {$b$};
  \node[active node] (c) at (4, 0) {$c$};
  \node[sink node] (t) at (6, 0) {$t$};

  \node[above, font=\footnotesize] at (a.north) {$e(a) = 3$};
  \node[below, font=\footnotesize] at (b.south) {$e(b) = 0$};
  \node[above, font=\footnotesize] at (c.north) {$e(c) = 1$};

  \draw[flow edge] (s) -- node[above left, font=\footnotesize] {5} (a);
  \draw[flow edge] (s) -- node[below left, font=\footnotesize] {2} (b);
  \draw[flow edge] (a) -- node[above, font=\footnotesize, pos=0.35] {2} (c);
  \draw[flow edge] (b) -- node[below, font=\footnotesize, pos=0.35] {2} (c);
  \draw[flow edge] (c) -- node[above, font=\footnotesize] {3} (t);
\end{tikzpicture}
\caption{Pré-fluxo com relaxação da conservação: os vértices intermediários $a$ e $c$ possuem excesso positivo ($e > 0$), sendo classificados como ativos.}
\label{fig:prefluxo}
\end{figure}


\section{Invariantes de Altura e Mecânica de Push-Relabel}\label{sec:push_relabel}

A execução de \textit{Push-Relabel}~\cite{goldberg1988new} consiste em aplicar duas operações sobre vértices ativos:

\begin{itemize}
    \item \textbf{Push:} Aplicado sobre um arco residual admissível $(u,v)$, isto é, com $c_f(u,v) > 0$ e $h(u) = h(v) + 1$. A quantidade enviada é $\delta = \min(e(u), c_f(u,v))$. Se $\delta = c_f(u,v)$, o push é saturante; caso contrário, é não saturante (zerando $e(u)$)~\cite{cook1998combinatorial}.

    \item \textbf{Relabel:} Aplicado quando o vértice ativo $u$ não possui arcos admissíveis ($c_f(u,v) > 0 \implies h(u) \le h(v)$). A altura é atualizada para $h(u) = 1 + \min \{h(v) \mid c_f(u,v) > 0\}$, restaurando a existência de arcos admissíveis.
\end{itemize}

Com seleção de vértices ativos em fila FIFO (\textit{First-In, First-Out}), o número de pushes não saturantes é amortizado, resultando em complexidade de tempo $\mathcal{O}(|V|^3)$~\cite{goldberg1988new}. A implementação está detalhada no Código~\ref{lst:push_relabel} (Apêndice~\ref{ap:push_relabel}).

\begin{figure}[!ht]
\centering
\begin{tikzpicture}[thick,node distance=2.2cm]

\tikzsubcaption{(1.5,4.0)}{Antes de \texttt{push}: arco admissível}

\node[active node] (u) at (0.0,3.0) {$u$};
\node[font=\footnotesize,above=2pt of u] {$h{=}3,\,e{=}5$};

\node[flow node,fill=blue!10] (v) at (3.0,3.0) {$v$};
\node[font=\footnotesize,above=2pt of v] {$h{=}2,\,e{=}1$};

\draw[->,line width=1.5pt,green!60!black] (u) -- node[above,font=\footnotesize] {$c_f=4$} (v);

\tikzsubcaption{(1.5,1.0)}{Depois de \texttt{push(u, v, $\delta$=4)}}
\node[active node] (u2) at (0.0,0.0) {$u$};
\node[font=\footnotesize,above=2pt of u2] {$h{=}3,\,e{=}1$};
\node[flow node,fill=blue!10] (v2) at (3.0,0.0) {$v$};
\node[font=\footnotesize,above=2pt of v2] {$h{=}2,\,e{=}5$};
\draw[nontree edge] (u2) -- node[above,font=\footnotesize,gray] {$c_f=0$} (v2);

\tikzsubcaption{(8.5,4.0)}{Antes de \texttt{relabel}: vértice $u$ sem arco admissível}
\node[gap node,fill=orange!30] (r1) at (7.0,3.0) {$u$};
\node[font=\footnotesize,above=2pt of r1] {$h{=}2,\,e{=}3$};
\node[flow node,fill=gray!10] (r2) at (10.0,3.0) {$w$};
\node[font=\footnotesize,above=2pt of r2] {$h{=}3$};
\draw[nontree edge] (r1) -- node[above,font=\footnotesize,align=center] {$c_f=2$\\inadmissível} (r2);

\tikzsubcaption{(8.5,1.0)}{Depois de \texttt{relabel(u)}: $h(u)\leftarrow 4$}
\node[gap node,fill=orange!50] (r3) at (7.0,0.0) {$u$};
\node[font=\footnotesize,above=2pt of r3] {$h{=}4,\,e{=}3$};
\node[flow node,fill=gray!10] (r4) at (10.0,0.0) {$w$};
\node[font=\footnotesize,above=2pt of r4] {$h{=}3$};
\draw[->,green!60!black,line width=1.5pt] (r3) -- node[above,font=\footnotesize] {admissível} (r4);
\end{tikzpicture}
\caption{Operações elementares de Push-Relabel. Esquerda: Push admissível de $u$ para $v$ ($c_f > 0$ e $h(u)=h(v)+1$). Direita: Relabel de $u$ ajustando a altura para criar novo arco admissível.}
\label{fig:push_relabel}
\end{figure}


\section{Cuidados na Inicialização e Aritmética de Infinitos}

Na inicialização, satura-se diretamente cada arco de saída da fonte $s$ (\lstinline{push_flow(edge_id, cap)}), sem alocar uma variável de excesso para a fonte. Em aplicações onde capacidades assumem valor \lstinline{INF} (como na segmentação de imagens), subtrair \lstinline{INF} de uma variável de excesso limitaria a fonte aritmeticamente (\lstinline{INF - INF = 0}), esgotando indevidamente a capacidade disponível. A injeção direta evita o rastreamento desnecessário de excesso em $s$.


\section{Otimização: A Heurística de Lacuna (\textit{Gap Heuristic})}\label{sec:push_relabel_gap}

Sem heurísticas, um vértice pode ser rerrotulado até $2|V|-1$ vezes, gerando $\mathcal{O}(|V|^2)$ relabels. A Heurística de Lacuna (\textit{Gap Heuristic}) detecta descontinuidades na distribuição de alturas por meio de um vetor de contagem (\lstinline{height_count}).

Se após um relabel o número de vértices no nível $k$ ($0 < k < |V|$) torna-se zero, estabelece-se uma lacuna (\textit{gap}). Como a invariante $h(u) \le h(v) + 1$ impede arcos residuais de decrescerem mais de um nível por passo, qualquer vértice com $h(u) > k$ não possui caminho residual até o sumidouro ($h(t) = 0$). A Figura~\ref{fig:gap_prova} ilustra a separação no digrafo residual.

\begin{figure}[!ht]
\centering
\begin{tikzpicture}[thick]
  \draw[layer line] (-2, 0) -- (6, 0) node[right,black,font=\footnotesize] {$h=0$ (sumidouro $t$)};
  \foreach \y/\h in {1.5/1, 3.0/2, 6.0/4} {
    \draw[layer line] (-2, \y) -- (6, \y) node[right,black,font=\footnotesize] {$h=\h$};
  }
  \draw[red!40,dashed,line width=1.5pt] (-2, 4.5) -- (6, 4.5) node[right,red!80!black,font=\footnotesize] {$h=3$ (\textbf{LACUNA})};

  \node[source node,font=\footnotesize] (t) at (2, 0) {$t$};
  \node[flow node,fill=yellow!20,font=\footnotesize] (a) at (2, 1.5) {$a$};
  \node[flow node,fill=yellow!20,font=\footnotesize] (b) at (0.5, 3.0) {$b$};
  \node[flow node,fill=yellow!20,font=\footnotesize] (c) at (3.5, 3.0) {$c$};

  \node[gap node,fill=red!20,font=\footnotesize] (u) at (2, 6.0) {$u$};

  \draw[->,blue!80!black] (b) -- (a);
  \draw[->,blue!80!black] (c) -- (a);
  \draw[->,blue!80!black] (a) -- (t);

  \draw[->,red,line width=1.5pt] (u) -- node[left,pos=0.25,font=\footnotesize,text=black,align=center,xshift=-2mm,fill=white,inner sep=2pt] {Queda de 2 níveis\\(inadmissível)} (b);
  \draw[->,red,line width=1.5pt] (u) -- node[right,pos=0.25,font=\footnotesize,text=black,align=center,xshift=2mm,fill=white,inner sep=2pt] {$h(u) \not\le h(c) + 1$\\(inadmissível)} (c);

\end{tikzpicture}
\caption{Inalcançabilidade na Heurística de Lacuna: a condição $h(u) \le h(v) + 1$ restringe o avanço a no máximo um nível por arco residual. A ausência de vértices no nível $h=3$ impede a conexão residual entre vértices de nível $h \ge 4$ e o sumidouro.}
\label{fig:gap_prova}
\end{figure}

Ao detectar a lacuna em $k$, o algoritmo ajusta a altura de todos os vértices com $h(u) > k$ para $|V| + 1$~\cite{korte2018combinatorial}, eliminando relabels desnecessários e direcionando o fluxo restante para a fonte.

\begin{figure}[!ht]
\centering
\begin{tikzpicture}[thick]

\begin{scope}[xshift=0cm]
\tikzsubcaption{(1.5,6.8)}{(a) Lacuna detectada no nível $k=3$}
\foreach \h/\label/\fill in {
  0/h=0/green!20, 1.2/h=1/yellow!20, 2.4/h=2/yellow!20,
  4.8/h=4/red!20, 6.0/h=5/red!20} {
  \node[circle,draw,fill=\fill,minimum size=8mm,font=\footnotesize] at (1.5,\h) {\label};
}
\node[font=\footnotesize,red!60!black] at (1.5,3.6) {\textbf{lacuna}};
\draw[<->,red,dashed] (0.2,4.2) -- node[left,font=\footnotesize] {nív. 3 vazio} (0.2,3.0);
\end{scope}

\node at (5.0,3.0) {\Huge $\Rightarrow$};

\begin{scope}[xshift=7cm]
\tikzsubcaption{(1.5,6.8)}{(b) Vértices com $h>k$ elevados a $|V|+1$}
\foreach \h/\label/\fill in {
  0/h=0/green!20, 1.2/h=1/yellow!20, 2.4/h=2/yellow!20} {
  \node[circle,draw,fill=\fill,minimum size=8mm,font=\footnotesize] at (1.5,\h) {\label};
}
\node[gap node,font=\footnotesize] (top) at (1.5,6.0) {$|V|{+}1$};
\draw[->,orange!80!black,line width=1.5pt] (1.5,3.3) -- node[right,font=\footnotesize,orange!80!black] {atualização} (1.5,5.05);
\end{scope}
\end{tikzpicture}
\caption{Heurística de Lacuna: (a) o nível $k=3$ esvazia-se. (b) Vértices com $h > k$ recebem altura $|V|+1$, acelerando a devolução de excessos.}
\label{fig:gap_heuristica}
\end{figure}


\section{Política de Seleção: FIFO versus Maior Altura}

A implementação com fila FIFO (Código~\ref{lst:push_relabel}, Apêndice~\ref{ap:push_relabel}) atinge complexidade $\mathcal{O}(|V|^3)$.

A versão com política de \textbf{Maior Altura} (\textit{Highest Label}) (Código~\ref{lst:push_relabel_improved}, Apêndice~\ref{ap:push_relabel_improved}) seleciona a cada iteração o vértice ativo de maior valor $h(u)$.

O gerenciamento de vértices ativos é feito por um vetor de listas indexado por altura (\lstinline{active_buckets}), com uma variável \lstinline{highest_active}:
\begin{itemize}
    \item \textbf{Inserção:} Quando $v$ torna-se ativo, é adicionado a \lstinline{active_buckets[h(v)]}, atualizando \lstinline{highest_active} se necessário ($\mathcal{O}(1)$).
    \item \textbf{Extração:} Remove-se de \lstinline{active_buckets[highest_active]}. Se a lista esvaziar, decrementa-se o índice ($\mathcal{O}(1)$ amortizado).
    \item \textbf{Relabel:} Atualizada a altura de $u$, ele é transferido para o balde da nova altura.
\end{itemize}

A política de Maior Altura associada à Heurística de Lacuna assegura a complexidade $\mathcal{O}(|V|^2 \sqrt{|A|})$~\cite{goldberg1988new}.

\begin{table}[!ht]
\centering
\small
\setlength{\tabcolsep}{5pt}
\begin{tabular}{p{4.2cm} l p{7.5cm}}
\toprule
\multicolumn{1}{c}{\textbf{Algoritmo}} & \multicolumn{1}{c}{\textbf{Complexidade}} & \multicolumn{1}{c}{\textbf{Mecanismo / Estratégia}} \\
\midrule
\multicolumn{3}{c}{\bannertext{\hyperref[sec:fluxo_primal]{Vertente Primal: Caminhos Aumentantes ($e(v)=0$)}}} \\
\midrule
\hyperref[ap:ford_fulkerson]{Ford-Fulkerson} & $\mathcal{O}(|A| \cdot |f^*|)$ & DFS (busca caminhos aumentantes arbitrários) \\
\hyperref[ap:edmonds_karp]{Edmonds-Karp} & $\mathcal{O}(|V| \cdot |A|^2)$ & BFS (caminhos de menor cardinalidade de arcos) \\
\hyperref[ap:dinic]{Dinic} & $\mathcal{O}(|V|^2 \cdot |A|)$ & Digrafo de níveis e fluxos bloqueadores \\
\midrule
\multicolumn{3}{c}{\bannertext{\hyperref[sec:prefluxo]{Vertente Dual: Pré-fluxo e Alturas (Invariante de Corte)}}} \\
\midrule
\hyperref[ap:push_relabel]{Push-Relabel (FIFO)} & $\mathcal{O}(|V|^3)$ & Fila FIFO de vértices ativos \\
\hyperref[ap:push_relabel_improved]{Push-Relabel (Gap)} & $\mathcal{O}(|V|^2 \sqrt{|A|})$ & Seleção por maior altura e \textit{Gap Heuristic} \\
\bottomrule\end{tabular}
\caption{Complexidades assintóticas e estratégias dos algoritmos de fluxo máximo implementados.}
\label{tab:max_flow}
\end{table}


\section{Síntese Comparativa: Primalidade e Dualidade no Fluxo Máximo}\label{sec:sintese_max_flow}

A Tabela~\ref{tab:max_flow} sintetiza as propriedades das abordagens implementadas para Fluxo Máximo:
\begin{itemize}
    \item \textbf{Vertente Primal (Caminhos Aumentantes):} Mantém a conservação de fluxo e busca caminhos de $s$ a $t$ em $D_f$. O algoritmo de Dinic apresenta melhor desempenho assintótico ao consolidar múltiplos aumentos por fase via fluxos bloqueadores no Digrafo de Níveis.
    \item \textbf{Vertente Dual (Pré-fluxo e Alturas):} Relaxa a conservação mantendo a ausência de caminhos aumentantes. As operações são locais e orientadas por alturas. A política de maior altura com Gap Heuristic atinge $\mathcal{O}(|V|^2 \sqrt{|A|})$, oferecendo menor tempo de CPU em grafos densos.
\end{itemize}

O algoritmo de Dinic destaca-se em redes esparsas com capacidades unitárias, enquanto Push-Relabel com Gap Heuristic sobressai em redes densas de alta capacidade.

